% One flexible-data frame with the rate switch, against one classic frame of
% the same payload.
% Everything here is in millimetres, because \framefield draws at y = 0 in
% millimetres and mixing it with centimetre coordinates puts the row labels in
% a different place from their own rows.
\begin{tikzpicture}[font=\sffamily\scriptsize]

% ================================================================ classic
\begin{scope}
  \node[signame] at (-2mm,3mm) {classic, 8 bytes};
  \framefield{0}{16}{fid}{arbitration\\{\tiny identifier and control}}
  \framefield{16}{40}{fdata}{data, 8 bytes}
  \framefield{56}{20}{fcrc}{checksum\\{\tiny and acknowledge}}
  \node[chlabel, anchor=west] at (80mm,3mm) {one rate throughout};
\end{scope}

% ================================================================ flexible-data
\begin{scope}[yshift=-22mm]
  \node[signame] at (-2mm,3mm) {flexible-data, 8 bytes, rate switch set};
  \framefield{0}{16}{fid}{arbitration\\{\tiny identifier and control}}
  \framefield{16}{12}{fdata}{data}
  \framefield{28}{16}{fcrc}{checksum\\{\tiny stronger}}
  \node[chlabel, anchor=west] at (48mm,3mm) {the same payload, in less time};
  \draw[faultedge, <->] (16mm,-3mm) -- (28mm,-3mm);
  \node[tlabel, anchor=north, text=red!60!black] at (22mm,-3.5mm)
    {the data phase, four times faster};
\end{scope}

% ================================================================ the marks
\draw[busstub] (0mm,-32mm) -- (0mm,10mm);
\node[tlabel, anchor=north, text width=32mm, align=center] at (0mm,-34mm)
  {the start-of-frame bit: the controller captures its timestamp here, in
   hardware, before anything else happens};

\draw[busstub] (16mm,-32mm) -- (16mm,-14mm);
\node[tlabel, anchor=north, text width=28mm, align=center] at (34mm,-34mm)
  {the rate switches here, and only here};

% ================================================================ the readings
\node[umlnote, text width=54mm, anchor=north west] at (0mm,-48mm)
  {Arbitration runs at the nominal rate in both frames, because every node on
   the bus has to take part in it. Only the data phase changes speed, and only
   between the two marked points. That is why a one megabit transceiver still
   gives the full payload and the stronger checksum: it simply never switches.};

\node[umlnote, text width=54mm, anchor=north west] at (60mm,-48mm)
  {The measurement in step 7 uses this structure: with the timestamp counter
   clocked from the bit time, the difference between two captured timestamps is
   a frame's duration in bit times, which confirms both the frame structure and
   the ratio between the phases with nothing attached to the board.};
\end{tikzpicture}
